%! Composition of Functions — Unit 6, Lesson 6.6 — Homework (Answer Key)
\documentclass[10pt]{article}
\usepackage{algebra2-article}
\usepackage{algebra2-key}   % pulls in -boxes; do NOT also load -boxes
\newcommand{\TallMath}[1]{$\displaystyle #1\rule[-1.4em]{0pt}{3.2em}$}
% In-formula answer slots (\ans is text-mode and must never appear inside $...$).
\newcommand{\mans}[1]{{\color{keyred}\mathbf{#1}}}
\newcommand{\mexp}[1]{{\color{keyred}#1}}
\newcommand{\lgrid}[4]{%
  \draw[step=1, linegray!40, very thin] (#1,#3) grid (#2,#4);
  \draw[->, semithick, charcoal] (#1,0)--(#2,0) node[below right, font=\scriptsize]{$x$};
  \draw[->, semithick, charcoal] (0,#3)--(0,#4) node[left, font=\scriptsize]{$y$};}

\begin{document}

\pageheader{Unit 6, Lesson 6.6}{Homework --- Answer Key}
\namedateperiod

\vspace{0.10in}

\begin{notesbox}{Practice}
Work \textbf{inside out}, and write the middle number every time. An answer with no middle step is
an answer you cannot check.

\begin{enumerate}[leftmargin=1.9em, itemsep=11pt]
  \item \textbf{From rules.} Let $f(x)=x^{2}-3$ and $g(x)=2x+5$.
    \begin{center}
    \renewcommand{\arraystretch}{1.8}
    \begin{tabularx}{\linewidth}{@{}X X X@{}}
    (a) \; $(f\circ g)(1)=\mans{46}$ &
    (b) \; $(g\circ f)(1)=\mans{1}$ &
    (c) \; $(f\circ f)(2)=\mans{-2}$ \\
    (d) \; $(g\circ g)(-1)=\mans{11}$ &
    (e) \; $(f\circ g)(-2)=\mans{-2}$ &
    (f) \; $(g\circ f)(-3)=\mans{17}$ \\
    \end{tabularx}
    \end{center}
    {\small \emph{Middles:} (a) $g(1)=7$ \; (b) $f(1)=-2$ \; (c) $f(2)=1$ \; (d) $g(-1)=3$ \;
    (e) $g(-2)=1$ \; (f) $f(-3)=6$.}

    Parts (a) and (b) use the same functions and the same input. Whose fault is the difference?
    \ans{The order --- (a) runs $g$ first, (b) runs $f$ first.}

  \item \textbf{From a table.} Use only the values below.
    \begin{center}
    \renewcommand{\arraystretch}{1.4}
    \begin{tabular}{@{}|c|c|c|c|c|c|c|@{}}
    \hline
    $x$ & $-2$ & $-1$ & $0$ & $1$ & $2$ & $3$ \\
    \hline
    $p(x)$ & $3$ & $1$ & $0$ & $-1$ & $2$ & $-2$ \\
    \hline
    $q(x)$ & $-1$ & $2$ & $3$ & $0$ & $1$ & $-2$ \\
    \hline
    \end{tabular}
    \end{center}
    \begin{center}
    \renewcommand{\arraystretch}{1.7}
    \begin{tabularx}{\linewidth}{@{}X X X@{}}
    (a) \; $(p\circ q)(0)=\mans{-2}$ &
    (b) \; $(q\circ p)(0)=\mans{3}$ &
    (c) \; $(p\circ q)(-1)=\mans{2}$ \\
    (d) \; $(q\circ p)(2)=\mans{1}$ &
    (e) \; $(p\circ p)(1)=\mans{1}$ &
    (f) \; $(q\circ q)(3)=\mans{-1}$ \\
    \end{tabularx}
    \end{center}
    (g) Find the value of $x$ for which $(p\circ q)(x)=0$. \; $x=\mans{1}$ \\[2pt]
    \hspace*{0.3in}\emph{Hint: work backwards --- which input does $p$ send to $0$?}
    {\small \ans{$p(0)=0$, so we need $q(x)=0$, which happens at $x=1$.}}

  \item \textbf{Algebraically, with domains.} Let $f(x)=x-6$, \; $g(x)=\sqrt{x}$, \; and
        $h(x)=x^{2}+2$.
    \begin{center}
    \renewcommand{\arraystretch}{1.8}
    \begin{tabularx}{\linewidth}{@{}l X X@{}}
    \hline
    \textbf{Composition} & \textbf{Simplified rule} & \textbf{Domain} \\
    \hline
    (a) \; $(g\circ f)(x)$ & \ans{$\sqrt{x-6}$} & \ans{$x\ge6$} \\
    (b) \; $(f\circ g)(x)$ & \ans{$\sqrt{x}-6$} & \ans{$x\ge0$} \\
    (c) \; $(h\circ f)(x)$ & \ans{$(x-6)^{2}+2$} & \ans{all reals} \\
    (d) \; $(f\circ h)(x)$ & \ans{$x^{2}-4$} & \ans{all reals} \\
    (e) \; $(g\circ h)(x)$ & \ans{$\sqrt{x^{2}+2}$} & \ans{all reals} \\
    \hline
    \end{tabularx}
    \end{center}
    \vspace{0.55in}
    Two of these five have restricted domains and three do not. What do the restricted ones have in
    common that (e) --- which also contains a square root --- does not? \ansline{In (a) and (b) the
    square root can be handed a negative number, so a gate closes. In (e) the radicand $x^{2}+2$ is
    at least $2$ for every real $x$, so the square root is never in danger --- the presence of a
    radical does not restrict a domain; a radicand that can go negative does.}

  \item \textbf{From two graphs.} Below are $y=u(x)$ and $y=v(x)$.
    \begin{center}
    \begin{minipage}[c]{0.47\linewidth}
    \centering
    \begin{tikzpicture}[scale=0.34]
      \lgrid{-5}{5}{-5}{5}
      \draw[very thick, forest] (-4,2)--(0,-2)--(4,2);
      \fill[forest] (0,-2) circle (3.4pt);
      \node[forest, font=\scriptsize] at (-3.0,3.4) {$y=u(x)$};
    \end{tikzpicture}
    \end{minipage}
    \hfill
    \begin{minipage}[c]{0.47\linewidth}
    \centering
    \begin{tikzpicture}[scale=0.34]
      \lgrid{-5}{5}{-5}{5}
      \draw[very thick, navy] (-4,-5)--(4,3);
      \fill[navy] (1,0) circle (3.4pt);
      \node[navy, font=\scriptsize] at (2.6,3.6) {$y=v(x)$};
    \end{tikzpicture}
    \end{minipage}
    \end{center}
    \begin{center}
    \renewcommand{\arraystretch}{1.7}
    \begin{tabularx}{\linewidth}{@{}X X X@{}}
    (a) \; $(u\circ v)(0)=\mans{-1}$ &
    (b) \; $(v\circ u)(0)=\mans{-3}$ &
    (c) \; $(u\circ v)(-3)=\mans{2}$ \\
    (d) \; $(v\circ u)(-3)=\mans{0}$ &
    (e) \; $(u\circ v)(3)=\mans{0}$ &
    (f) \; $(v\circ u)(3)=\mans{0}$ \\
    \end{tabularx}
    \end{center}
    (g) Parts (e) and (f) came out \emph{equal}, even though (a) and (b) did not. Try $x=4$ as
    well: $(u\circ v)(4)=\mans{1}$ and $(v\circ u)(4)=\mans{1}$. Where on the number line
    do the two orders seem to agree, and what is special about $u$ there? \ansline{They agree for
    $x\ge1$. There $u$ is on its right-hand arm, where it is simply ``subtract $2$'' with no corner
    to fold the sign, so both orders reduce to $x-3$. To the left of the corner $u$ turns around and
    the two orders come apart.}

  \item \textbf{Error analysis.} Two classmates are asked for $(f\circ g)(x)$ where $f(x)=x^{2}$ and
        $g(x)=x+3$.
    \begin{center}
    \begin{tcolorbox}[colback=white, colframe=charcoal!45, boxrule=0.7pt, arc=1.5mm,
        width=0.72\linewidth, left=3mm, right=3mm, top=1.5mm, bottom=1.5mm]
    \centering\small
    \renewcommand{\arraystretch}{1.4}
    \begin{tabular}{@{}l l@{}}
    \textbf{Amara writes:} & $(f\circ g)(x)=x^{2}\cdot(x+3)=x^{3}+3x^{2}$ \\
    \textbf{Ben writes:} & $(f\circ g)(x)=x^{2}+3$ \\
    \end{tabular}
    \end{tcolorbox}
    \end{center}
    \vspace{-4pt}
    (a) What is the correct answer? \ans{$(x+3)^{2}=x^{2}+6x+9$} \\[3pt]
    (b) Evaluate all three expressions at $x=2$: correct $=\mans{25}$, Amara $=\mans{20}$,
    Ben $=\mans{7}$ \\[3pt]
    (c) Name each mistake in a few words --- they are two \emph{different} misunderstandings, not
    the same slip twice. \\[2pt]
    Amara: \ans{read $\circ$ as multiplication} \\[3pt]
    Ben: \ans{right idea, wrong order --- that is $(g\circ f)(x)$}

  \item \textbf{Explain.} Let $f(x)=x^{2}$ and $g(x)=\sqrt{x}$. The composition
        $(f\circ g)(x)=\left(\sqrt{x}\,\right)^{2}$ simplifies all the way down to just $x$ --- and
        yet its domain is \emph{not} all real numbers. Explain why, using the words \textbf{stage}
        (or \emph{gate}) and \textbf{domain}. Then say what $(g\circ f)(x)$ simplifies to instead,
        and why that one \emph{does} accept every real number. \ansline{Stage 1 is $g(x)=\sqrt{x}$,
        whose domain is $x\ge0$, so a negative input is turned away before stage 2 ever runs ---
        there is no $\sqrt{-9}$ to square. Simplifying to $x$ removes the radical from the
        \emph{formula} but cannot re-open a gate that already closed, so the domain of $f\circ g$ is
        $x\ge0$. \\[3pt]
        In the other order, $(g\circ f)(x)=\sqrt{x^{2}}=|x|$. Now stage 1 is squaring, which accepts
        every real number and always produces a non-negative result, so stage 2's gate is never
        tested --- the domain is all reals. The price is that the answer is $|x|$, not $x$: at
        $x=-9$ it returns $9$.}
\end{enumerate}
\end{notesbox}

\begin{extensionbox}
\textbf{The view from the mast.} In Lesson 6.4 you met the horizon formula: from an eye height of
$h$ feet above the water, you can see about
\[
  d(h)=\sqrt{1.5h} \ \text{ miles.}
\]
A sailor starts with their eyes $6$ feet above the deck and climbs the mast at a steady $6$ feet per
second, so after $t$ seconds their eye height is
\[
  h(t)=6t+6 \ \text{ feet.}
\]
\begin{enumerate}[label=(\alph*), leftmargin=1.8em, itemsep=5pt]
  \item Which function has to run first if you want the horizon distance after $t$ seconds ---
        $d$ or $h$? Write the composition you need in $\circ$ notation.
        \ans{$h$ runs first (time $\rightarrow$ height $\rightarrow$ distance): $(d\circ h)(t)$.}
  \item Build the rule. Show that $(d\circ h)(t)=\sqrt{1.5(6t+6)}$ simplifies to
        $3\sqrt{t+1}$. \ans{$1.5(6t+6)=9t+9=9(t+1)$, and $\sqrt{9(t+1)}=3\sqrt{t+1}$.}
  \item Complete the table. Every entry is a whole number of miles.
    \begin{center}
    \renewcommand{\arraystretch}{1.5}
    \begin{tabular}{@{}|c|c|c|c|c|@{}}
    \hline
    $t$ (seconds) & $0$ & $3$ & $8$ & $15$ \\
    \hline
    $(d\circ h)(t)$ (miles) & \ans{$3$} & \ans{$6$} & \ans{$9$} & \ans{$12$} \\
    \hline
    \end{tabular}
    \end{center}
  \item Your composition is $3\sqrt{t+1}$ --- which, in the language of Lesson 6.4, is the parent
        $y=\sqrt{t}$ with a horizontal shift and a vertical stretch. Name both, and explain what
        each one \emph{means} about the sailor. \ans{Shift \emph{left} $1$ and stretch vertically by
        $3$. The shift is the $6$-foot head start: at $t=0$ the sailor is already as high as one
        second of climbing would have taken them, so the graph begins one unit into the parent
        rather than at its endpoint. The stretch is the $\sqrt{1.5\cdot6}=3$ buried in the
        formula --- every horizon distance is three times the parent's.}
  \item It takes $3$ seconds of climbing to go from $3$ miles of visibility to $6$. How long to get
        from $6$ miles to $12$? Which feature of a square root graph made that answer inevitable?
        \ans{$3\sqrt{t+1}=12\Rightarrow t=15$, so $12$ more seconds --- four times as long.
        Doubling the output requires \emph{quadrupling} $t+1$, because the graph flattens: the same
        gain in height buys less and less new horizon.}
  \item The algebra is happy to compute $(d\circ h)(-1)=0$. What does the \emph{situation} say about
        that input, and what is the domain of the composition \emph{in context}?
        \ans{$t=-1$ is one second before the climb began, which is not a time this model describes
        (and the sailor is never at eye height $0$ anyway). In context the domain is $t\ge0$, so the
        usable graph starts at $(0,3)$ --- three miles, not zero.}
\end{enumerate}
\end{extensionbox}

\begin{spiralbox}
\textbf{Coming next --- Lesson 6.7: Inverse Functions.} A few of today's pairs did something the
others did not: composed in \emph{either} order, they handed back the $x$ you started with. That is
not a coincidence and it is not a curiosity --- it is the definition of a pair of \textbf{inverse}
functions, and $f(g(x))=g(f(x))=x$ is the test. In 6.7 you will learn to \emph{build} the partner
for a given function (swap and solve), see the pair as reflections of each other over the line
$y=x$, and watch the domain and range trade places. And the loose thread from today's Section 6
finally gets tied: $\left(\sqrt{x}\,\right)^{2}$ gives back $x$ only when $x\ge0$, which is exactly
why $y=x^{2}$ has to have its domain restricted before it is allowed a square-root inverse --- the
question you first raised all the way back in Lesson 6.0.
\end{spiralbox}

\begin{teachernote}
\small Items 1--2 are fluency and should take five minutes; items 3--6 are the assessed content.
\begin{itemize}[leftmargin=1.3em, nosep]
  \item \textbf{Item 2(g) is the first backwards question} students have seen on compositions, and
        it is the seed of 6.7: to hit a target output you have to un-do the outer machine first.
        Expect it to be skipped; it is worth a minute of class time.
  \item \textbf{Item 3's follow-up is the diagnostic.} The wrong generalization ---
        ``a square root means a restricted domain'' --- passes rows (a) and (b) and fails row (e).
        A student who marks (e) as $x\ge0$ or similar has memorized a symbol instead of a mechanism.
  \item \textbf{Item 4(g)} rewards a group that noticed the corner. The honest answer is $x\ge1$:
        past the corner of $u$, both orders are ``subtract $3$.'' Accept any correct description of
        that region; the reasoning matters more than the interval notation.
  \item \textbf{Item 5 separates the two failure modes.} Amara's is a notation error ($\circ$ read
        as $\cdot$), Ben's is an order error. They need different corrections, and the $x=2$ test
        settles both without argument.
  \item \textbf{Item 6 is the bridge to 6.7 --- read these before planning tomorrow.} A student who
        can write it clearly is ready for the restricted-domain argument; a student who writes
        ``because you can't square root a negative'' without connecting it to \emph{which stage} is
        not, and 6.7 should open by re-running Notes Section 6 for them.
  \item \textbf{The extension} deliberately ends where Lesson 6.4 began: a real formula turns out to
        be nothing but the parent with a shift and a stretch --- and now students can see \emph{why},
        because the shift and the stretch are the two machines they composed.
\end{itemize}
\end{teachernote}

\end{document}
